% ==========================================================================
%  Guia Completo Unificado — Prova Única 2026
%  Macroeconomia I (PPGEco-UFES, 2026)
%  Prof. Ricardo Ramalhete Moreira
%  Cobre: Solow (cap.1), RCK + Diamond (cap.2), RBC (cap.5),
%         Rigidez Nominal (cap.6), NK DSGE / Fischer / Taylor / Calvo / CGG (cap.7)
%  Lista I (Q1-Q8) + Lista II (Q1-Q11)
% ==========================================================================
\documentclass[12pt, a4paper]{article}

% --------------------------------------------------------------------------
% Packages
% --------------------------------------------------------------------------
\usepackage[utf8]{inputenc}
\usepackage[T1]{fontenc}
\usepackage{lmodern}
\usepackage[brazil]{babel}
\usepackage{microtype}

\usepackage[top=2.5cm, bottom=2.5cm, left=3cm, right=3cm, headheight=15pt]{geometry}

\usepackage{amsmath, amssymb, amsthm}
\usepackage{mathtools}
\usepackage{bm}

\usepackage{booktabs}
\usepackage{array}
\usepackage{multicol}
\usepackage{xcolor}
\usepackage{hyperref}
\usepackage{enumitem}
\usepackage{tcolorbox}
\tcbuselibrary{breakable}
\usepackage{fancyhdr}
\usepackage{titlesec}

% --------------------------------------------------------------------------
% Cores
% --------------------------------------------------------------------------
\definecolor{cyanrbc}{HTML}{3b9dbd}
\definecolor{orangerbc}{HTML}{d97b39}
\definecolor{grayrbc}{HTML}{6b7280}
\definecolor{greenbox}{HTML}{166534}
\definecolor{greenbg}{HTML}{dcfce7}
\definecolor{boxbg}{HTML}{f5f0e8}
\definecolor{boxborder}{HTML}{3b9dbd}
\definecolor{provabg}{HTML}{fff7ed}
\definecolor{provaborder}{HTML}{d97b39}

% --------------------------------------------------------------------------
% Caixas (standard skin — sem skins/TikZ para evitar manchas negras)
% --------------------------------------------------------------------------
\newtcolorbox{resultbox}[1][]{standard, colback=boxbg, colframe=boxborder,
  boxrule=1pt, arc=3pt, left=6pt, right=6pt, top=4pt, bottom=4pt,
  fonttitle=\bfseries\small, #1}
\newtcolorbox{provabox}[1][]{standard, colback=provabg, colframe=provaborder,
  boxrule=1.2pt, arc=3pt, left=6pt, right=6pt, top=4pt, bottom=4pt,
  fonttitle=\bfseries\small, #1}
\newtcolorbox{formulabox}[1][]{standard, colback=greenbg!30, colframe=greenbox!60,
  boxrule=0.8pt, arc=3pt, left=6pt, right=6pt, top=4pt, bottom=4pt,
  fonttitle=\bfseries\small, #1}

% --------------------------------------------------------------------------
% Comandos matemáticos
% --------------------------------------------------------------------------
\newcommand{\E}{\mathbb{E}}
\newcommand{\dif}{\mathrm{d}}
\newcommand{\khat}{\hat{k}}
\newcommand{\chat}{\hat{c}}
\newcommand{\yhat}{\hat{y}}
\newcommand{\zhat}{\hat{z}}
\newcommand{\rstar}{r^*}
\newcommand{\kstar}{k^*}
\newcommand{\cstar}{c^*}
\newcommand{\pibar}{\bar{\pi}}

% --------------------------------------------------------------------------
% Espaçamento e parágrafos
% --------------------------------------------------------------------------
\setlength{\parskip}{0.5\baselineskip}
\setlength{\parindent}{0pt}

% --------------------------------------------------------------------------
% Títulos de seções
% --------------------------------------------------------------------------
\titleformat{\section}{\large\bfseries\color{cyanrbc}}{}{0em}{}[\vspace{-0.3em}\rule{\linewidth}{0.4pt}]
\titleformat{\subsection}{\normalsize\bfseries\color{grayrbc}}{}{0em}{}

% --------------------------------------------------------------------------
% Cabeçalho/rodapé
% --------------------------------------------------------------------------
\pagestyle{fancy}
\fancyhf{}
\lhead{\footnotesize\textcolor{grayrbc}{Guia Completo da Prova --- Macroeconomia I (PPGEco-UFES, 2026)}}
\rhead{\footnotesize\textcolor{grayrbc}{Solow $\cdot$ RCK $\cdot$ Diamond $\cdot$ RBC $\cdot$ Rigidez Nominal $\cdot$ NK DSGE}}
\cfoot{\small\thepage}
\renewcommand{\headrulewidth}{0.4pt}

% ==========================================================================
\begin{document}
% ==========================================================================

% --------------------------------------------------------------------------
% Capa
% --------------------------------------------------------------------------
\begin{titlepage}
\begin{center}
\vspace*{1.5cm}
{\Large\bfseries Universidade Federal do Espírito Santo}\\[0.5em]
{\large Programa de Pós-Graduação em Economia (PPGEco-UFES)}\\[2em]

\rule{\linewidth}{1.2pt}\\[0.8em]
{\LARGE\bfseries\color{cyanrbc} Guia Completo Unificado da Prova\\[0.4em]
\large Macroeconomia I --- 2026}\\[0.4em]
{\normalsize\color{grayrbc} Romer caps.\ 1, 2, 5, 6 e 7 (Aulas 1--9)}\\[0.8em]
\rule{\linewidth}{1.2pt}

\vspace{1em}
{\large Prof.\ Ricardo Ramalhete Moreira}

\vspace{1.5em}
\begin{tcolorbox}[standard, colback=provabg, colframe=provaborder, width=0.92\textwidth,
  arc=4pt, boxrule=1.2pt, left=8pt, right=8pt, top=6pt, bottom=6pt]
\centering
\textbf{\textcolor{provaborder}{Prova Única:}} Lista I (Q1--Q8) e Lista II (Q1--Q11) são para a
\textbf{MESMA} prova. Este guia unifica todo o conteúdo:\\[0.4em]
\textbf{Solow} (cap.\ 1) $\cdot$ \textbf{RCK + Diamond} (cap.\ 2) $\cdot$ \textbf{RBC} (cap.\ 5)
$\cdot$ \textbf{Rigidez Nominal} (cap.\ 6) $\cdot$ \textbf{NK DSGE / Fischer / Taylor / Calvo / CGG} (cap.\ 7)
\end{tcolorbox}

\vspace{1.2em}
\begin{tabular}{lll}
\toprule
\textbf{Parte} & \textbf{Conteúdo} & \textbf{Lista} \\
\midrule
I & Solow: crescimento exógeno, Regra de Ouro & --- \\
II & Diamond OLG: ineficiência dinâmica & --- \\
III & RCK: bem-estar, Keynes-Ramsey, diagrama de fases & Lista I Q1--Q3 \\
IV & RBC: capital, trabalho, PTF & Lista I Q4--Q8 \\
V & Modelos de Ajuste de Preços (Fischer/Taylor/Calvo/CEE) & --- \\
VI & Modelo NK Canônico CGG + Soluções & Lista II Q1--Q11 \\
Ap. A & Questões-tipo com Fischer/Taylor/Calvo/CGG & --- \\
Ap. B & Tabelas de calibração unificadas & --- \\
\bottomrule
\end{tabular}

\vspace{1.5em}
\begin{tcolorbox}[standard, colback=boxbg, colframe=boxborder, width=0.92\textwidth,
  arc=4pt, boxrule=1pt, left=8pt, right=8pt, top=6pt, bottom=6pt]
\centering
\textbf{Referências principais:}\\[0.3em]
D.\ Romer, \textit{Advanced Macroeconomics}, 5.\textordfeminine\ ed., McGraw-Hill (2019). Caps.\ 1, 2, 5, 6, 7.\\[0.2em]
J.\ Galí, \textit{Monetary Policy, Inflation, and the Business Cycle}, Princeton (2008).\\[0.2em]
Slides das Aulas 8--9 (Cap.\ 7): Fischer (1977), Taylor (1980), Calvo (1983), CGG (2000).
\end{tcolorbox}

\vfill
{\small Guia unificado para a prova única de Macroeconomia I --- PPGEco-UFES 2026.}
\end{center}
\end{titlepage}

% --------------------------------------------------------------------------
% Formulário de Referência Rápida (pág. 2) — 3 colunas
% --------------------------------------------------------------------------
\newpage
\thispagestyle{fancy}

\begin{center}
{\large\bfseries\color{cyanrbc} Formulário de Referência Rápida --- Prova Única 2026}\\[0.3em]
\rule{\linewidth}{0.6pt}
\end{center}

\vspace{0.3em}
\begin{multicols}{3}

\textbf{\textcolor{cyanrbc}{SOLOW + RCK}}

\medskip
\textbf{Eq.\ fundamental Solow:}
\[
\dot{k} = sf(k) - (n+g+\delta)k
\]

\textbf{EE Cobb-Douglas:}
\[
k^* = \!\left(\frac{s}{n+g+\delta}\right)^{\!1/(1-\alpha)}
\]

\textbf{Regra de Ouro:}
\[
f'(k_{gold}) = n+g+\delta,\quad s_{gold} = \alpha
\]

\textbf{Keynes-Ramsey:}
\[
\frac{\dot{c}}{c} = \frac{f'(k)-\delta-\rho-\theta g}{\theta}
\]

\textbf{EE RCK:}
\[
r^* = \rho + \theta g,\quad \beta(r^*+1-\delta)=1
\]

\textbf{Convergência RCK:}
\[
\rho - n - (1-\theta)g > 0
\]

\textbf{Crescimento balanceado:}
\[
g_K = g_Y = n+g,\quad g_{Y/L} = g
\]

\textbf{Resíduo de Solow:}
\[
g_{TFP} = g_Y - \alpha g_K - (1-\alpha)g_L
\]

\columnbreak

\textbf{\textcolor{cyanrbc}{DIAMOND + RBC}}

\medskip
\textbf{Diamond --- poupança:}
\[
s_t = \frac{\beta}{1+\beta}\,w_t
\]

\textbf{Diamond --- mapa:}
\[
k_{t+1} = \frac{\beta(1-\alpha)}{(1+\beta)(1+n)}k_t^\alpha
\]

\textbf{Inef.\ dinâmica:} $r^* < n$ $\Rightarrow$ $k^* > k_{gold}$

\textbf{RBC --- Euler capital:}
\[
\beta(r^*+1-\delta)=1,\quad I^*=\delta K^*
\]

\textbf{DML e Lazer:}
\[
\text{DML} = \frac{b}{1-N_t},\quad l^* = \frac{bC}{w}
\]

\textbf{Subs.\ intertemporal:}
\[
\frac{1-l_1}{1-l_2} = \beta(1+r)\frac{w_2}{w_1}
\]

\textbf{Calibração:}
\[
b = \frac{w^* l^*}{C^*} \approx 1{,}55
\]

\textbf{3 canais PTF:}
(1) direto $\cdot$ (2) investimento $\cdot$ (3) suavização

\columnbreak

\textbf{\textcolor{cyanrbc}{NK MODELOS (Cap.\ 7)}}

\medskip
\textbf{IS NK:}
\[
\ln Y_t = \ln Y_{t+1} - \frac{1}{\theta}r_t
\]

\textbf{$\phi$ (rigidez real):}
\[
\phi = \theta+\gamma-1,\quad p_t^* = \phi m_t + (1-\phi)p_t
\]

\textbf{CPNK Calvo:}
\[
\pi_t = \kappa y_t + \beta E_t\pi_{t+1}
\]
\[
\kappa = \frac{\alpha[1-(1-\alpha)\beta]\phi}{1-\alpha}
\]

\textbf{CPNK Híbrida CEE:}
\[
\pi_t = \gamma\pi_{t-1} + (1-\gamma)E_t\pi_{t+1} + \chi y_t
\]

\textbf{CGG (Taylor forward-looking):}
\[
r_t = \phi_\pi E_t[\pi_{t+1}] + \phi_y E_t[y_{t+1}] + u_t^{MP}
\]

\textbf{Fischer eq.\ (20):}
\[
y_t = \tfrac{1}{1+\phi}[E_{t-1}m_t - E_{t-2}m_t] + [m_t - E_{t-1}m_t]
\]

\textbf{Taylor eq.\ (22):}
\[
y_t = \lambda y_{t-1} + \tfrac{1+\lambda}{2}u_t
\]

\textbf{Princípio de Taylor:} $\phi_\pi > 1$

\end{multicols}

\vspace{0.5em}
\begin{provabox}[title={Alertas Críticos da Prova}]
\textbf{(1)} Convergência RCK: $\rho - n \mathbf{-} (1-\theta)g > 0$ (sinal $\mathbf{-}$, não $+$).
\textbf{(2)} CPNK do professor: $\kappa = \alpha[1-(1-\alpha)\beta]\phi/(1-\alpha)$ (notação das slides), \textbf{diferente} de Galí: $\kappa = (1-\alpha)(1-\alpha\beta)\phi/\alpha$ (ambas equivalentes, mas notação diverge na prova).
\textbf{(3)} Regra CGG: \textbf{forward-looking} $E_t[\pi_{t+1}]$, não contemporânea $\pi_t$.
\textbf{(4)} Fischer/Taylor: $\phi$ mede rigidez real; $\phi<1$ $\Rightarrow$ maior rigidez real $\Rightarrow$ maiores efeitos reais da moeda.
\end{provabox}

% --------------------------------------------------------------------------
% Sumário
% --------------------------------------------------------------------------
\newpage
\tableofcontents
\newpage

% ==========================================================================
\section*{\large\color{cyanrbc} Parte I --- Modelo de Solow: Crescimento Exógeno (Cap.\ 1)}
\addcontentsline{toc}{section}{Parte I --- Modelo de Solow (Cap.\ 1)}
% ==========================================================================

\subsection*{I.1 Forma Intensiva e Equação Fundamental}

\textbf{Ponto de partida.} A função de produção agregada é $Y = F(K, AL)$, onde $A$ cresce à taxa $g$ e $L$ cresce à taxa $n$. Com Retornos Constantes de Escala (RCE): $y \equiv Y/(AL) = f(k)$, $k \equiv K/(AL)$.

\textbf{Passo 1 --- dinâmica de $k$.} Diferenciando $k = K/(AL)$:
\[
\dot{k} = \frac{\dot{K}}{AL} - (g+n)k.
\]

\textbf{Passo 2 --- lei de movimento do capital.} $\dot{K} = sY - \delta K$ $\Rightarrow$ $\dot{K}/(AL) = sf(k) - \delta k$. Substituindo:

\begin{formulabox}[title={Equação Fundamental de Solow}]
\[
\boxed{\dot{k} = sf(k) - (n+g+\delta)k}
\]
$sf(k)$ = investimento; $(n+g+\delta)k$ = investimento de manutenção.
\end{formulabox}

\subsection*{I.2 Estado Estacionário e Cobb-Douglas}

No EE: $sf(k^*) = (n+g+\delta)k^*$. Com $f(k) = k^\alpha$:

\begin{resultbox}[title={Estado Estacionário de Solow --- Cobb-Douglas}]
\[
k^* = \left(\frac{s}{n+g+\delta}\right)^{\!1/(1-\alpha)}, \quad
y^* = \left(\frac{s}{n+g+\delta}\right)^{\!\alpha/(1-\alpha)}, \quad
c^* = (1-s)y^*.
\]
\end{resultbox}

\subsection*{I.3 Regra de Ouro}

Maximizar $c^* = f(k^*) - (n+g+\delta)k^*$ em $k^*$:

\textbf{FOC:} $f'(k_{gold}) = n+g+\delta$.

Com Cobb-Douglas: $\alpha k_{gold}^{\alpha-1} = n+g+\delta$ $\Rightarrow$ $k_{gold} = (\alpha/(n+g+\delta))^{1/(1-\alpha)}$.

\textbf{Taxa de poupança áurea:} $s_{gold} = f'(k_{gold})\cdot k_{gold}/f(k_{gold}) = \alpha$.

\begin{resultbox}[title={Regra de Ouro de Solow}]
\[
f'(k_{gold}) = n+g+\delta, \qquad s_{gold} = \alpha.
\]
Se $s > s_{gold}$: sobre-poupança (ineficiência dinâmica --- reduzir $s$ Pareto-melhora).
\end{resultbox}

\subsection*{I.4 Crescimento Balanceado e Contabilidade}

No EE, $k = K/(AL)$ constante implica $g_K = n+g$. Com Cobb-Douglas: $g_Y = n+g$ e $g_{Y/L} = g$.

\begin{resultbox}[title={Contabilidade do Crescimento (Resíduo de Solow)}]
\[
g_{TFP} \equiv g_Y - \alpha g_K - (1-\alpha)g_L.
\]
\end{resultbox}

\subsection*{I.5 Convergência Condicional}

Linearizando ao redor de $k^*$: taxa de convergência $|\mu| = (1-\alpha)(n+g+\delta) > 0$.

\begin{provabox}[title={Solow vs.\ RCK --- Diferença Fundamental}]
\textbf{Solow:} $s$ exógena. \textbf{RCK:} $s$ endógena (otimização). No RCK: $r^* = \rho+\theta g > n+g$ (retorno áureo), logo $k_{RCK}^* < k_{gold}$.
\end{provabox}

% ==========================================================================
\section*{\large\color{cyanrbc} Parte II --- Modelo de Diamond: Gerações Sobrepostas (Cap.\ 2)}
\addcontentsline{toc}{section}{Parte II --- Modelo de Diamond (Cap.\ 2)}
% ==========================================================================

\subsection*{II.1 Configuração e Regra de Poupança}

Dois períodos de vida. Utilidade: $U = \ln(c_{1t}) + \beta\ln(c_{2,t+1})$.

\textbf{FOC:} $-1/(w_t - s_t) + \beta/s_t = 0$ $\Rightarrow$ $\beta(w_t-s_t) = s_t$.

\begin{resultbox}[title={Poupança Ótima no Modelo de Diamond}]
\[
s_t = \frac{\beta}{1+\beta}\,w_t.
\]
A poupança é \emph{independente} da taxa de juros com log-utilidade.
\end{resultbox}

\subsection*{II.2 Mapa Dinâmico de Capital}

\[
k_{t+1} = \frac{s_t}{1+n} = \frac{\beta(1-\alpha)}{(1+\beta)(1+n)}\,k_t^\alpha \equiv \Phi\, k_t^\alpha.
\]

\begin{resultbox}[title={Mapa Dinâmico de Diamond}]
\[
k_{t+1} = \Phi\, k_t^\alpha, \qquad \Phi \equiv \frac{\beta(1-\alpha)}{(1+\beta)(1+n)},\qquad k^{D*} = \Phi^{1/(1-\alpha)}.
\]
\end{resultbox}

\subsection*{II.3 Regra de Ouro e Ineficiência Dinâmica}

Com $\delta=1$: Regra de Ouro exige $r_{gold} = n$. Se $k^{D*} > k_{gold}$: $r^* < n$ --- acumulação excessiva de capital.

\begin{provabox}[title={Ineficiência Dinâmica no Diamond vs.\ RCK}]
\textbf{Diamond:} cada geração não internaliza efeitos intergeracionais $\Rightarrow$ possível $k^{D*} > k_{gold}$ (ineficiência dinâmica).\\
\textbf{RCK:} família com horizonte infinito internaliza tudo $\Rightarrow$ equilíbrio Pareto-ótimo.
\end{provabox}

% ==========================================================================
\section*{\large\color{cyanrbc} Parte III --- Modelo RCK: Questões 1--3 da Lista I (Cap.\ 2)}
\addcontentsline{toc}{section}{Parte III --- Modelo RCK: Questões 1--3 (Lista I)}
% ==========================================================================

\setcounter{section}{0}

\section{Efeitos sobre o Bem-Estar: $\theta$, $\rho$, $n$, $g$ (Q1)}

\subsection*{1.1 Função Utilidade, $\beta$ e Bem-estar}

A função utilidade da família, reescrita em termos de consumo por trabalho efetivo $c(t)$:
\[
U = B\int_0^\infty e^{-\beta t}\!\left[\frac{c(t)^{1-\theta}}{1-\theta}\right]dt, \qquad \boxed{\beta \equiv \rho - n - (1-\theta)g.}
\]
$B = A(0)^{1-\theta}L(0)/H > 0$ é uma constante. $\beta > 0$ garante convergência da integral (condição de transversalidade).

\subsection*{1.2 Estado Estacionário no RCK}

\textbf{$\dot{c}=0$:} $f'(k^*) = \delta+\rho+\theta g$ $\Rightarrow$ $r^* = \rho+\theta g$ $\Rightarrow$ $k^* = (\alpha/(\rho+\theta g+\delta))^{1/(1-\alpha)}$.

\textbf{$\dot{k}=0$:} $c^* = f(k^*) - (\delta+n+g)k^*$.

\subsection*{1.3 Estática Comparativa}

\begin{resultbox}[title={Resumo da Estática Comparativa (Q1)}]
\begin{center}
\begin{tabular}{lccccp{5cm}}
\toprule
Parâmetro & $r^*$ & $k^*$ & $c^*$ & $W^*$ & Mecanismo \\
\midrule
$\theta\uparrow$ & $\uparrow$ & $\downarrow$ & $\downarrow$ & $\downarrow$ & Menor crescimento do consumo ($\dot{c}/c \downarrow$) \\
$\rho\uparrow$   & $\uparrow$ & $\downarrow$ & $\downarrow$ & $\downarrow\downarrow$ & Impaciência + $\beta\uparrow$ (duplo) \\
$n\uparrow$      & $-$        & $-$          & $\downarrow$ & $\uparrow$ & Mais produto total; maior consumo por família \\
$g\uparrow$      & $\uparrow$ & $\downarrow$ & ambíguo      & $\uparrow$ & Produtividade $\uparrow$; $\beta\downarrow$ (menos desconto) \\
\bottomrule
\end{tabular}
\end{center}
\end{resultbox}

\begin{resultbox}[title={Gabarito Oficial Q1 --- Respostas (a)--(d)}]
\textbf{(a) $\theta\uparrow$:} maior aversão ao risco reduz o crescimento intertemporal do consumo ($\dot{C}/C = (r-\rho)/\theta$). Maximização da utilidade em \textbf{nível menor de bem-estar} da família.\\[3pt]
\textbf{(b) $\rho\uparrow$:} famílias valorizam mais a utilidade corrente em relação à futura $\Rightarrow$ menor taxa de crescimento do consumo $\Rightarrow$ \textbf{menor bem-estar}.\\[3pt]
\textbf{(c) $n\uparrow$:} RCK é modelo de pleno emprego. Mais trabalhadores $\Rightarrow$ produto total maior $\Rightarrow$ maior consumo por família $\Rightarrow$ \textbf{maior bem-estar} (veja o papel de $n$ em $\beta = \rho-n-(1-\theta)g$: $\beta\downarrow$).\\[3pt]
\textbf{(d) $g\uparrow$:} economia eleva produto per capita via produtividade. Na função utilidade, $\beta = \rho - n - (1-\theta)g$ decresce $\Rightarrow$ menor desconto sobre valores futuros $\Rightarrow$ \textbf{melhoria do bem-estar}.
\end{resultbox}

\begin{provabox}[title={Q1 --- Atenção em (c) e (d)}]
Armadilha frequente: confundir $c^*$ per efficiency unit (que cai com $n\uparrow$) com o bem-estar familiar total (que sobe, via $\beta\downarrow$ e via mais membros consumindo). Para $g\uparrow$: o efeito sobre $k^*$ é negativo ($r^*\uparrow$), mas o efeito sobre $W^*$ é positivo ($\beta\downarrow$ e produtividade maior).
\end{provabox}

% --------------------------------------------------------------------------
\section{Taxa de Crescimento Ótima do Consumo (Q2)}
% --------------------------------------------------------------------------

\subsection*{2.1 Método do Lagrangiano (abordagem oficial)}

Função utilidade e restrição orçamentária (com $r$ constante e riqueza $W^*$ dada):
\[
U = \int_0^\infty e^{-\rho t}\!\left[\frac{C(t)^{1-\theta}}{1-\theta}\right]\frac{L(t)}{H}dt, \qquad
\text{s.a.}\quad \int_0^\infty e^{-rt}C(t)\frac{L(t)}{H}dt = W^*.
\]

Lagrangiano:
\[
\mathcal{L} = \int_0^\infty e^{-\rho t}\!\left[\frac{C(t)^{1-\theta}}{1-\theta}\right]\frac{L(t)}{H}dt + \lambda\!\left[W^* - \int_0^\infty e^{-rt}C(t)\frac{L(t)}{H}dt\right].
\]

\textbf{CPO em $C(t)$:}
\[
e^{-\rho t}C(t)^{-\theta}\frac{L(t)}{H} - \lambda\, e^{-rt}\frac{L(t)}{H} = 0 \;\Rightarrow\; e^{-\rho t}C(t)^{-\theta} = \lambda\, e^{-rt}. \tag{CPO}
\]

\textbf{Diferenciando ambos os lados em relação a $t$:}
\[
-\rho\, e^{-\rho t}C^{-\theta} + e^{-\rho t}(-\theta C^{-\theta-1}\dot{C}) = -r\lambda\, e^{-rt}.
\]
Usando (CPO) para substituir $\lambda e^{-rt} = e^{-\rho t}C^{-\theta}$ e cancelando $e^{-\rho t}C^{-\theta}$:
\[
-\theta\frac{\dot{C}}{C} - \rho + r = 0.
\]

\begin{formulabox}[title={Equação de Euler do Consumo (Q2 --- Método Lagrangiano)}]
\[
\boxed{\frac{\dot{C}(t)}{C(t)} = \frac{r - \rho}{\theta}.}
\]
Em termos de consumo por trabalho efetivo $\dot{c} = \dot{C}/\dot{A}$:
\[
\frac{\dot{c}(t)}{c(t)} = \frac{\dot{C}}{C} - \frac{\dot{A}}{A} = \frac{r-\rho}{\theta} - g.
\]
\end{formulabox}

\subsection*{2.2 Hamiltoniano de Valor Corrente (método alternativo)}

\[
\mathcal{H} = \frac{c^{1-\theta}-1}{1-\theta} + \lambda\!\left[f(k) - c - (n+g+\delta)k\right].
\]

\textbf{FOC em $c$:} $c^{-\theta} = \lambda$. \textbf{Equação coadjunta:} $\dot{\lambda} = \rho\lambda - \lambda[f'(k)-(n+g+\delta)]$.

Derivando $\lambda = c^{-\theta}$: $\dot{\lambda}/\lambda = -\theta\dot{c}/c = \rho + \delta - f'(k)$.

\begin{formulabox}[title={Regra de Keynes-Ramsey (forma geral)}]
\[
\boxed{\frac{\dot{c}}{c} = \frac{f'(k) - \delta - \rho - \theta g}{\theta}}
\]
No EE ($f'(k^*) = \delta + \rho + \theta g$): $\dot{c}^* = 0$. Com $r = f'(k)-\delta$: recupera-se $\dot{C}/C = (r-\rho)/\theta$.
\end{formulabox}

\subsection*{2.3 Condição de Convergência}

Ao longo do BGP, o integrando cresce como $e^{[(1-\theta)g - (\rho-n)]t}$. Para convergir:

\begin{formulabox}[title={Condição de Convergência/Transversalidade (RCK)}]
\[
\boxed{\rho - n - (1-\theta)g > 0}
\]
\end{formulabox}

\begin{provabox}[title={Q2 --- O Sinal que Todo Mundo Erra}]
A condição correta é $\rho - n \mathbf{-} (1-\theta)g > 0$ (sinal $\mathbf{-}$ antes de $(1-\theta)g$). A alternativa errada usa $\rho - n + (1-\theta)g > 0$ (sinal $+$) --- erro mais cobrado em provas anteriores.
\end{provabox}

% --------------------------------------------------------------------------
\section{Diagrama de Fases: Regiões e Dinâmica (Q3)}
% --------------------------------------------------------------------------

\subsection*{3.1 As Duas Isóclinas}

\textbf{$\dot{c}=0$:} $k = k^*$ (vertical). À esquerda: $\dot{c}>0$; à direita: $\dot{c}<0$.

\textbf{$\dot{k}=0$:} $c = f(k)-(\delta+n+g)k$ (côncava). Acima: $\dot{k}<0$; abaixo: $\dot{k}>0$.

\begin{resultbox}[title={Mapa das Quatro Regiões do Diagrama de Fases (Q3)}]
\begin{center}
\begin{tabular}{lllll}
\toprule
Região & Posição & $\dot{c}$ & $\dot{k}$ & Direção \\
\midrule
NW & Esq.\ $\dot{c}=0$, acima $\dot{k}=0$ & $>0$ & $<0$ & $\uparrow\leftarrow$ \\
SW & Esq.\ $\dot{c}=0$, abaixo $\dot{k}=0$ & $>0$ & $>0$ & $\uparrow\rightarrow$ \\
NE & Dir.\ $\dot{c}=0$, acima $\dot{k}=0$ & $<0$ & $<0$ & $\downarrow\leftarrow$ \\
SE & Dir.\ $\dot{c}=0$, abaixo $\dot{k}=0$ & $<0$ & $>0$ & $\downarrow\rightarrow$ \\
\bottomrule
\end{tabular}
\end{center}
NW (resposta da Q3): $\dot{c}>0$ e $\dot{k}<0$ --- trajetória \textbf{diverge} do EE.
\end{resultbox}

\subsection*{3.2 Estabilidade de Sela}

O Jacobiano no EE tem $\det(J^*) = (c^*/\theta)f''(k^*) < 0$: autovalores com sinais opostos $\Rightarrow$ \textbf{ponto de sela}. Existe uma única trajetória convergente (sela ótima).

% ==========================================================================
\section*{\large\color{cyanrbc} Parte IV --- Modelo RBC: Questões 4--8 da Lista I (Cap.\ 5)}
\addcontentsline{toc}{section}{Parte IV --- Modelo RBC: Questões 4--8 (Lista I)}
% ==========================================================================

% --------------------------------------------------------------------------
\section{Investimento de Equilíbrio e Estabilidade do Capital (Q4)}
% --------------------------------------------------------------------------

\subsection*{4.1 Estado Estacionário do Capital no RBC}

\[
K_{t+1} = (1-\delta)K_t + I_t \;\Rightarrow\; \text{No EE: } I^* = \delta K^*.
\]

\begin{resultbox}[title={Investimento e Euler no EE do RBC}]
\[
I^* = \delta K^*, \qquad \beta(r^* + 1 - \delta) = 1 \;\Rightarrow\; r^* = \frac{1-\beta(1-\delta)}{\beta}.
\]
Para $\beta=0{,}99$, $\delta=0{,}025$: $r^* \approx 3{,}5\%$ por período.
\end{resultbox}

\subsection*{4.2 Log-linearização e Estabilidade ($|A|<1$)}

Após log-linearização: $\hat{k}_{t+1} = A\hat{k}_t + B\hat{z}_t$. Com calibração padrão ($\alpha=0{,}33$, $\beta=0{,}99$, $\delta=0{,}025$): $A \approx 0{,}962 < 1$. Meia-vida $\approx 17{,}9$ trimestres.

\begin{provabox}[title={Q4 --- Erros Frequentes}]
(1) $I^* = \delta K^*$ (correto), não $I^* > \delta K^*$. (2) Euler: $\beta(r^*+1-\delta)=1$ --- não $\beta(r^*-1+\delta)=1$. (3) $|A|<1$ necessário para estabilidade.
\end{provabox}

% --------------------------------------------------------------------------
\section{Desutilidade Marginal da Oferta de Trabalho (Q5)}
% --------------------------------------------------------------------------

Com $u(C_t, l_t) = \ln C_t + b\ln(1-N_t)$:

\begin{resultbox}[title={Desutilidade Marginal do Trabalho (DML)}]
\[
\text{DML} = \left|\frac{\partial u}{\partial N_t}\right| = \frac{b}{1-N_t}.
\]
DML cresce com $N_t$ (cada hora adicional custa mais) e com $b$ (maior preferência por lazer).
\end{resultbox}

\textbf{Especificação alternativa} $U = C_t - (1/\eta)N_t^\eta$: FOC intratemporal $N_t^{\eta-1} = W_t/P_t$ $\Rightarrow$ $N_t^* = (W_t/P_t)^{1/(\eta-1)}$, elasticidade $= 1/(\eta-1)$.

% --------------------------------------------------------------------------
\section{Lazer Ótimo: Efeitos de $r$ e $w$ Futuros (Q6)}
% --------------------------------------------------------------------------

\subsection*{6.1 Ótimo Intratemporal}

$bC_t/l_t = w_t$ $\Rightarrow$ $l_t^* = bC_t/w_t$.

\begin{resultbox}[title={Efeitos sobre Oferta de Trabalho (Q6)}]
$r_t\uparrow$: $C_t\downarrow$ (Euler) $\Rightarrow$ $l_t^*\downarrow$ $\Rightarrow$ $N_t\uparrow$ (substituição intertemporal de lazer).\\
$w_{t+1}\uparrow$: riqueza$\uparrow$ $\Rightarrow$ $C_t\uparrow$ $\Rightarrow$ $l_t^*\uparrow$ $\Rightarrow$ $N_t\downarrow$ (efeito riqueza domina).
\end{resultbox}

\subsection*{6.2 Fórmula da Substituição Intertemporal}

Combinando ótimos intraperiodais com Euler ($C_{t+1} = \beta(1+r_t)C_t$):

\begin{formulabox}[title={Substituição Intertemporal de Lazer/Trabalho}]
\[
\frac{l_t}{l_{t+1}} = \frac{w_{t+1}}{\beta(1+r_t)w_t} \quad\Longleftrightarrow\quad \frac{1-l_1}{1-l_2} = \frac{\beta(1+r)w_2}{w_1}.
\]
$r\uparrow$ $\Rightarrow$ lado direito$\uparrow$ $\Rightarrow$ $N_1/N_2\uparrow$ $\Rightarrow$ mais trabalho hoje.
\end{formulabox}

% --------------------------------------------------------------------------
\section{Razão Consumo-Lazer e Calibração de $b$ (Q7)}
% --------------------------------------------------------------------------

\begin{resultbox}[title={Calibração do Parâmetro $b$ (Q7)}]
\[
b = \frac{w^*(1-N^*)}{C^*} \approx 1{,}55,
\]
usando $N^* = 1/3$, $\alpha=0{,}33$, $\beta=0{,}99$, $\delta=0{,}025$.
\end{resultbox}

\begin{center}
\begin{tabular}{lll}
\toprule
Quantidade EE & Fórmula & Valor \\
\midrule
$r^*$ & $1/\beta - (1-\delta)$ & $0{,}0351$ \\
$k^*$ & $(\alpha/(r^*+\delta))^{1/(1-\alpha)}$ & $\approx 11{,}57$ \\
$y^*$ & $(k^*)^\alpha$ & $\approx 2{,}41$ \\
$w^*$ & $(1-\alpha)y^*$ & $\approx 1{,}61$ \\
$C^*$ & $y^*-\delta k^*$ & $\approx 2{,}12$ \\
$b$ & $w^*(1-N^*)/C^*$ & $\approx 1{,}55$ \\
\bottomrule
\end{tabular}
\end{center}

% --------------------------------------------------------------------------
\section{Três Canais de Transmissão dos Choques de PTF (Q8)}
% --------------------------------------------------------------------------

A tecnologia segue AR(1): $\hat{z}_{t+1} = \rho_z\hat{z}_t + \varepsilon_{t+1}$, $\rho_z \approx 0{,}95$.

\begin{resultbox}[title={Quatro Canais de Transmissão do Choque PTF (Q8)}]
\begin{enumerate}[leftmargin=1.5em, itemsep=2pt]
\item \textbf{Direto}: $Z_t\uparrow$ $\Rightarrow$ $Y_t\uparrow$, $r_t\uparrow$, $w_t\uparrow$ ($\hat{Y}_0^{direto} = 1\%$).
\item \textbf{Investimento}: $r_t\uparrow$ $\Rightarrow$ $I_t\uparrow\uparrow$ $\Rightarrow$ $K_{t+s}\uparrow$ (persistência endógena). $\hat{I}_0 \approx 3{,}2\%$.
\item \textbf{Suavização}: $C_t$ sobe menos que $Y_t$ (agente distribui no tempo). $\hat{C}_0 \approx 0{,}33\%$.
\item \textbf{Trabalho}: $w_t\uparrow$, $r_t\uparrow$ $\Rightarrow$ $N_t\uparrow$ (substituição intertemporal) $\Rightarrow$ amplifica $Y_t$.
\end{enumerate}
Hierarquia: $\hat{I}_0 \approx 3{,}2\% > \hat{Y}_0 \approx 1\text{--}1{,}5\% > \hat{C}_0 \approx 0{,}33\%$.
\end{resultbox}

% ==========================================================================
\newpage
\section*{\large\color{cyanrbc} Parte V --- Modelos de Ajuste de Preços (Cap.\ 7, Aulas 8--9)}
\addcontentsline{toc}{section}{Parte V --- Modelos de Ajuste de Preços: Fischer/Taylor/Calvo/CEE}
% ==========================================================================

Esta parte apresenta o conteúdo das Aulas 8--9 (Romer Cap.\ 7), que o professor cobre com ênfase na sequência Fischer $\to$ Taylor $\to$ Calvo $\to$ CEE.

% ------------------------------------------------------------------
\subsection*{V.1 Framework Geral: Famílias, Firmas e IS NK}

\textbf{Preferências das famílias (Romer 2012, Cap.\ 7):}
\[
U = \sum_{t=0}^\infty \beta^t \left[U(C_t) - V(L_t)\right],
\]
com $U(C_t) = C_t^{1-\theta}/(1-\theta)$ (CRRA) e $V(L_t) = (B/\gamma)L_t^\gamma$, $B>0$, $\gamma>1$.

\textbf{FOC para oferta de trabalho:}
\[
V'(L_t) = U'(C_t)\cdot\frac{W_t}{P_t} \;\Longrightarrow\; BL_t^{\gamma-1} = C_t^{-\theta}\cdot\frac{W_t}{P_t}.
\]

Com $C_t = L_t = Y_t$ (produção um-a-um, todo produto consumido):
\begin{equation}
\frac{W_t}{P_t} = BY_t^{\theta+\gamma-1} = BY_t^{\phi}, \tag{eq.\ 6}
\end{equation}
onde definimos o \textbf{parâmetro de rigidez real}:
\[
\boxed{\phi \equiv \theta + \gamma - 1.}
\]
$\phi$ é o expoente que governa a inclinação do salário real em relação ao produto.

\textbf{IS Novo-Keynesiana (eq.\ 7):}
\begin{equation}
\ln Y_t = \ln Y_{t+1} - \frac{1}{\theta}r_t. \tag{eq.\ 7}
\end{equation}
Note que $1/\theta$ é a elasticidade de substituição intertemporal (EIS). Com $r_t$ acima do normal, $Y_t < Y_{t+1}$ (produto presente abaixo da trajetória futura).

\textbf{Firmas:} $Y_{it} = L_{it}$ (produção um-a-um). Demanda Dixit-Stiglitz: $Y_{it} = Y_t(P_{it}/P_t)^{-\eta}$.

\textbf{Lucro real da firma $i$ (eq.\ 8):}
\begin{equation}
R_t = Y_t\!\left[\left(\frac{P_{it}}{P_t}\right)^{1-\eta} - \frac{W_t}{P_t}\left(\frac{P_{it}}{P_t}\right)^{-\eta}\right]. \tag{eq.\ 8}
\end{equation}

\textbf{Problema de minimização (eq.\ 11):} Com incerteza, firmas minimizam desvio quadrático do preço ótimo:
\begin{equation}
\min_{p_i} \sum_t q_t\,(p_i - p_t^*)^2, \tag{eq.\ 11}
\end{equation}
onde $p_t^*$ é o preço ótimo se a firma pudesse reajustar livremente.

\textbf{Preço ótimo (eqs.\ 12--15):} Com $m_t = y_t + p_t$ (quantidade de moeda em logs):
\begin{equation}
p_t^* = \phi m_t + (1-\phi)p_t, \tag{eq.\ 14}
\end{equation}
e o preço escolhido pela firma:
\begin{equation}
p_i = \sum_t \omega_t E_0[\phi m_t + (1-\phi)p_t]. \tag{eq.\ 15}
\end{equation}

\begin{resultbox}[title={Intuição de $\phi$ como Rigidez Real}]
\textbf{Caso $\phi=1$:} $p_t^* = m_t$ --- o preço ótimo move-se um-a-um com a moeda. Sem rigidez real.\\
\textbf{Caso $\phi<1$:} $p_t^*$ depende também de $p_t$ (preço geral). Maior dependência de $p_t$ $\Rightarrow$ maior \textbf{complementaridade estratégica} $\Rightarrow$ firmas relutam em afastar seu preço do nível geral $\Rightarrow$ maior \textbf{rigidez real}.\\
\textbf{Conclusão:} $\phi$ menor = maior rigidez real = maiores efeitos reais da política monetária.
\end{resultbox}

% ------------------------------------------------------------------
\subsection*{V.2 Modelo de Fischer (Preços Predeterminados, Contratos Escalonados)}

\textbf{Referência:} Fischer (1977). Contratos de longa duração, preços \emph{predeterminados} (fixados antes do período, mas podem diferir entre coortes).

\textbf{Estrutura de escalonamento:}
\begin{itemize}[leftmargin=1.5em, itemsep=2pt]
  \item Metade das firmas fixa o preço no período anterior ($t-1$) para vigorar em $t$ e $t+1$.
  \item Metade das firmas fixa o preço dois períodos antes ($t-2$) para vigorar em $t-1$ e $t$.
\end{itemize}

\textbf{Nível de preços (eq.\ 16):}
\begin{equation}
p_t = \frac{1}{2}(p_t^1 + p_t^2), \tag{eq.\ 16}
\end{equation}
onde $p_t^1$ é o preço fixado em $t-1$ e $p_t^2$ o fixado em $t-2$.

\textbf{Derivação passo a passo de $p_t^1$ e $p_t^2$:}

\textit{Preço $p_t^1$} (fixado em $t-1$, minimiza $E_{t-1}[(p_t^1 - p_t^*)^2 + (p_{t+1}^1 - p_{t+1}^*)^2]$):
\[
p_t^1 = \frac{1}{2}(E_{t-1}[p_t^*] + E_{t-1}[p_{t+1}^*]).
\]
Usando $p_t^* = \phi m_t + (1-\phi)p_t$ e tomando expectativas em $t-1$:
\[
p_t^1 \approx E_{t-1}[p_t^*] = \phi E_{t-1}[m_t] + (1-\phi)E_{t-1}[p_t].
\]

\textit{Preço $p_t^2$} (fixado em $t-2$):
\[
p_t^2 = E_{t-2}[p_t^*] = \phi E_{t-2}[m_t] + (1-\phi)E_{t-2}[p_t].
\]

\textbf{Resultado para o nível de preços (eq.\ 19):}
\begin{equation}
p_t = E_{t-2}m_t + \frac{\phi}{1+\phi}\left[E_{t-1}m_t - E_{t-2}m_t\right]. \tag{eq.\ 19}
\end{equation}

\textbf{Produto (eq.\ 20):} Usando $y_t = m_t - p_t$ e a eq.\ (19):
\begin{equation}
\boxed{y_t = \frac{1}{1+\phi}\left[E_{t-1}m_t - E_{t-2}m_t\right] + \left[m_t - E_{t-1}m_t\right].} \tag{eq.\ 20}
\end{equation}

\begin{resultbox}[title={Interpretação do Modelo de Fischer (eqs.\ 19--20)}]
\textbf{Componente 1:} $[E_{t-1}m_t - E_{t-2}m_t]/(1+\phi)$ --- revisão de expectativas entre $t-2$ e $t-1$, amplificada pelo fator $1/(1+\phi)$.\\
\textbf{Componente 2:} $[m_t - E_{t-1}m_t]$ --- surpresa monetária em $t$ (inovação), tem efeito real de magnitude \textbf{1} (pleno).\\
\textbf{Rigidez real:} $\phi\downarrow$ $\Rightarrow$ $1/(1+\phi)\uparrow$ $\Rightarrow$ revisões de expectativas têm \textbf{maior efeito real}. Maior rigidez real ($\phi$ baixo) amplifica os efeitos reais de choques monetários antecipados.
\end{resultbox}

\begin{provabox}[title={Fischer vs.\ Neoclássico --- Para a Prova}]
No modelo neoclássico (preços plenamente flexíveis), $y_t = m_t - p_t = 0$ para qualquer $m_t$ antecipado (neutralidade da moeda). No Fischer, apenas a \textbf{surpresa} $[m_t - E_{t-1}m_t]$ tem efeito real pleno; a revisão $[E_{t-1}m_t - E_{t-2}m_t]$ tem efeito \emph{parcial} $1/(1+\phi) < 1$. Se $\phi=1$ (sem rigidez real), o componente de revisão tem efeito $1/2$ --- ainda não-neutro devido ao escalonamento.
\end{provabox}

% ------------------------------------------------------------------
\subsection*{V.3 Modelo de Taylor (Preços Fixos, Contratos Escalonados)}

\textbf{Referência:} Taylor (1980). Preços \emph{fixos} (não apenas predeterminados) por dois períodos. A diferença em relação ao Fischer é que o preço não pode ser reajustado entre os dois períodos do contrato.

\textbf{Processo da moeda (random walk):}
\[
m_t = m_{t-1} + u_t, \quad u_t \sim i.i.d.
\]

\textbf{Preço escolhido em $t$ (eq.\ 21):} Firma que fixa o preço em $t$ usa:
\begin{equation}
x_t = \frac{1}{2}(p_{it}^* + E_t p_{it+1}^*) = \frac{1}{2}\left\{[\phi m_t + (1-\phi)p_t] + [\phi E_t m_{t+1} + (1-\phi)E_t p_{t+1}]\right\}. \tag{eq.\ 21}
\end{equation}

\textbf{Nível de preços:} $p_t = (x_t + x_{t-1})/2$ (metade das firmas usa preço desta semana, metade da semana passada).

\textbf{Solução por coeficientes indeterminados:} Supor $x_t = x_{t-1} + \lambda u_t$ (solução linear em $u_t$). Substituindo em (21) e na lei de movimento do produto:

\begin{equation}
\boxed{y_t = \lambda y_{t-1} + \frac{1+\lambda}{2}u_t.} \tag{eq.\ 22}
\end{equation}

O coeficiente $\lambda \in [0,1)$ satisfaz:
\[
\lambda = \frac{1-\phi}{1+\phi}.
\]

\begin{resultbox}[title={Modelo de Taylor --- Resultados (eq.\ 22)}]
\textbf{Caso $\phi=1$:} $\lambda = 0$ $\Rightarrow$ $y_t = u_t/2$. Sem persistência do produto: cada choque tem efeito imediato mas não perdura. Equivalente a preços plenamente flexíveis (neutralidade da moeda para choques passados).\\[0.3em]
\textbf{Caso $\phi<1$ (rigidez real):} $\lambda = (1-\phi)/(1+\phi) > 0$ $\Rightarrow$ $y_t$ tem \textbf{persistência positiva}. O produto de hoje depende do produto de ontem. Choques monetários têm efeitos que se prolongam no tempo --- gradualmente convergindo para zero.\\[0.3em]
\textbf{Intuição:} Com $\phi<1$, firmas dependem fortemente de $p_t$ (nível geral) para fixar seu preço $x_t$. Como $p_t$ e $x_{t-1}$ estão correlacionados, $x_t$ herda a inércia de $x_{t-1}$, gerando persistência do produto.
\end{resultbox}

\begin{provabox}[title={Comparação Fischer e Taylor --- Para a Prova}]
\textbf{Fischer:} preços predeterminados (podem ser revistos entre contratos se houver nova informação disponível). O escalonamento gera não-neutralidade mesmo sem rigidez real ($\phi=1$).\\
\textbf{Taylor:} preços \emph{fixos} (não há revisão intra-contrato). Rigidez real ($\phi<1$) é necessária para gerar \emph{persistência} do produto (não apenas efeitos imediatos).\\
\textbf{Ambos:} $\phi$ pequeno $\Rightarrow$ maiores efeitos reais. A chave é a \emph{complementaridade estratégica}: firmas relutam em afastar seu preço do preço dos concorrentes.
\end{provabox}

% ------------------------------------------------------------------
\subsection*{V.4 Modelo de Calvo e CPNK}

\textbf{Referência:} Calvo (1983). Em cada período, cada firma pode reajustar o preço com probabilidade $\alpha$ (independente do estado e do histórico).

\textbf{Nível de preços (eq.\ 23):}
\begin{equation}
p_t = \alpha x_t + (1-\alpha)p_{t-1}, \tag{eq.\ 23}
\end{equation}
onde $x_t$ é o preço escolhido pelas firmas que re-otimizam.

\textbf{Inflação (eq.\ 24):}
\begin{equation}
\pi_t \equiv p_t - p_{t-1} = \alpha(x_t - p_{t-1}). \tag{eq.\ 24}
\end{equation}

\textbf{Derivação da CPNK (eq.\ 25):}

\textit{Passo 1.} As firmas que re-otimizam em $t$ escolhem $x_t$ que minimiza o desvio esperado de $p_\tau^*$ em todos os períodos futuros em que o preço ainda vigorará:
\[
x_t = (1-\alpha\beta)\sum_{s=0}^\infty (\alpha\beta)^s E_t[p_{t+s}^*].
\]

\textit{Passo 2.} Usando $p_t^* = \phi m_t + (1-\phi)p_t$ e $m_t = y_t + p_t$:
\[
p_t^* = \phi y_t + p_t.
\]

\textit{Passo 3.} Forma recursiva: $x_t = \phi y_t + p_t + \alpha\beta E_t[x_{t+1} - p_t]$.

\textit{Passo 4.} Usando $\pi_t = \alpha(x_t - p_{t-1})$, logo $x_t = p_{t-1} + \pi_t/\alpha$:

\textit{Passo 5.} Após manipulação algébrica:

\begin{formulabox}[title={Curva de Phillips NK de Calvo (eq.\ 25) --- Notação das Slides}]
\[
\boxed{\pi_t = \kappa\, y_t + \beta\, E_t\pi_{t+1},}
\]
\[
\boxed{\kappa \equiv \frac{\alpha[1-(1-\alpha)\beta]\phi}{1-\alpha}.}
\]
\textbf{Nota:} Esta é a fórmula do professor (slides). A fórmula de Galí usa notação diferente: $\kappa_{Galí} = (1-\alpha)(1-\alpha\beta)(\omega+\sigma^{-1})/\alpha$. Com $\omega = \sigma^{-1} = \phi/2$, as fórmulas são equivalentes.
\end{formulabox}

\textbf{Interpretação dos parâmetros de $\kappa$:}
\begin{itemize}[leftmargin=1.5em, itemsep=2pt]
  \item $\phi\uparrow$ $\Rightarrow$ $\kappa\uparrow$: menor rigidez real $\Rightarrow$ inflação mais responsiva ao produto.
  \item $\alpha\uparrow$ (mais firmas ajustam) $\Rightarrow$ $\kappa\uparrow$: mais firmas re-otimizando acelera o pass-through.
  \item $\beta\uparrow$ $\Rightarrow$ $\kappa\downarrow$: firmas mais pacientes ponderam mais o futuro, atenuando a pressão presente.
\end{itemize}

\begin{provabox}[title={CPNK Calvo: Forward-looking --- Para a Prova}]
A CPNK de Calvo é \textbf{puramente forward-looking}: $\pi_t$ depende de $E_t\pi_{t+1}$ (futuro), não de $\pi_{t-1}$ (passado). Implicação crucial: uma desinflação \emph{crível} ($E_t\pi_{t+1}\downarrow$) pode reduzir $\pi_t$ sem recessão ($y_t = 0$). Isso contradiz a evidência empírica de Ball (1994b) --- daí a necessidade do modelo CEE.
\end{provabox}

% ------------------------------------------------------------------
\subsection*{V.5 Inércia Inflacionária e Evidência Empírica}

\textbf{O problema do modelo puro de Calvo.} A CPNK $\pi_t = \kappa y_t + \beta E_t\pi_{t+1}$ implica que:
\begin{itemize}[leftmargin=1.5em, itemsep=2pt]
  \item Para reduzir $\pi_t$, basta reduzir $E_t\pi_{t+1}$ (credibilidade do BC).
  \item Uma desinflação crível é \textbf{sem custo}: $y_t$ pode permanecer em 0.
  \item Isso contradiz a evidência: desinflações historicamente custam produto.
\end{itemize}

\textbf{Evidência empírica --- Ball (1994b):}
\begin{itemize}[leftmargin=1.5em, itemsep=2pt]
  \item 28 episódios de desinflação em 9 países, período 1960--1990.
  \item Todos os episódios foram atribuídos à política monetária contracionista.
  \item \textbf{27 de 28 episódios}: produto ficou \emph{abaixo} do normal durante a desinflação.
  \item Conclusão: desinflação tem \textbf{custo real} --- a CPNK pura é inconsistente com os dados.
\end{itemize}

\textbf{``Inércia inflacionária'' $\neq$ inflação serialmente correlacionada.}
O termo ``inércia'' não significa apenas que $\pi_t$ é correlacionado com $\pi_{t-1}$, mas sim que \textbf{reduções na inflação têm custo real} (razão de sacrifício positiva). A CPNK pura gera correlação serial (pois $E_t\pi_{t+1}$ é persistente), mas sem custo de sacrifício.

\textbf{Solução --- CEE (Christiano, Eichenbaum e Evans, 2005):}

Modificação do Calvo: entre re-otimizações, as firmas indexam seu preço à inflação do período anterior. Isso gera inércia \emph{real}.

\textbf{Nível de preços com indexação (eq.\ 27):}
\begin{equation}
p_t = (1-\alpha)(p_{t-1} + \pi_{t-1}) + \alpha x_t. \tag{eq.\ 27}
\end{equation}

\textbf{CPNK Híbrida (eq.\ 28):}
\begin{equation}
\boxed{\pi_t = \gamma\pi_{t-1} + (1-\gamma)E_t\pi_{t+1} + \chi y_t,} \tag{eq.\ 28}
\end{equation}
com $0 \leq \gamma \leq 1$ e $\chi > 0$.

\textbf{Galí-Gertler (1999) --- CPNK Híbrida Estimada:}
\begin{equation}
\pi_t = \gamma_b\pi_{t-1} + \gamma_f E_t\pi_{t+1} + \kappa(y_t - \bar{y}_t) + e_t. \tag{eq.\ 26}
\end{equation}

\begin{resultbox}[title={CPNK Híbrida CEE vs.\ Calvo Puro}]
\begin{center}
\begin{tabular}{lll}
\toprule
Modelo & Equação & Custo de desinflação \\
\midrule
Calvo puro & $\pi_t = \kappa y_t + \beta E_t\pi_{t+1}$ & Zero (se crível) \\
CEE Híbrido & $\pi_t = \gamma\pi_{t-1} + (1-\gamma)E_t\pi_{t+1} + \chi y_t$ & Positivo se $\gamma > 1-\gamma$ \\
\bottomrule
\end{tabular}
\end{center}
\textbf{Condição para custo positivo:} $\gamma > 1-\gamma$ $\Leftrightarrow$ $\gamma > 1/2$.\\
Com $\gamma > 1/2$: para $E_t\pi_{t+1} < \pi_t$ (desinflação esperada), é necessário $y_t < 0$ para satisfazer a CPNK. Portanto, a desinflação exige recessão.
\end{resultbox}

\begin{provabox}[title={Inércia Inflacionária --- Para a Prova}]
\textbf{Armadilha frequente:} confundir ``inflação serialmente correlacionada'' com ``inércia inflacionária'' no sentido de Ball (1994b). A CPNK pura gera correlação serial (via $E_t\pi_{t+1}$ persistente), mas não razão de sacrifício positiva.\\
\textbf{Mecanismo CEE:} a indexação ``ancora'' o preço à inflação passada. Quando $\pi$ cai, as firmas que não re-otimizam ainda ajustam pela inflação passada (elevada), pressionando o nível de preços para cima. Para que $\pi$ caia, é necessário que $y < 0$ comprima os custos marginais --- daí o custo real.
\end{provabox}

% ==========================================================================
\newpage
\section*{\large\color{cyanrbc} Parte VI --- Modelo NK Canônico CGG + Questões Lista II (Q1--Q11)}
\addcontentsline{toc}{section}{Parte VI --- Modelo NK Canônico CGG + Lista II (Q1--Q11)}
% ==========================================================================

\subsection*{VI.1 Modelo CGG (Clarida-Galí-Gertler 2000)}

O modelo canônico Novo-Keynesiano consiste em três equações com choques AR(1):

\textbf{IS dinâmica:}
\begin{equation}
y_t = E_t[y_{t+1}] - \frac{1}{\theta}r_t + u_t^{IS}, \tag{IS}
\end{equation}

\textbf{CPNK (do professor):}
\begin{equation}
\pi_t = \beta E_t[\pi_{t+1}] + \kappa y_t + u_t^\pi, \quad 0<\beta<1,\; \kappa>0, \tag{CPNK}
\end{equation}

\textbf{Regra de Taylor CGG (forward-looking):}
\begin{equation}
\boxed{r_t = \phi_\pi E_t[\pi_{t+1}] + \phi_y E_t[y_{t+1}] + u_t^{MP}, \quad \phi_\pi>0,\; \phi_y\geq 0.} \tag{CGG}
\end{equation}

\textbf{Atenção:} O professor usa a regra \textbf{forward-looking} de Clarida-Galí-Gertler (2000) --- a taxa de juros reage às expectativas $E_t[\pi_{t+1}]$ e $E_t[y_{t+1}]$, não às variáveis correntes $\pi_t$ e $y_t$ (que é a especificação de Taylor 1993).

\textbf{Três choques AR(1):}
\[
u_t^{IS} = \rho_{IS}u_{t-1}^{IS} + e_t^{IS};\quad u_t^\pi = \rho_\pi u_{t-1}^\pi + e_t^\pi;\quad u_t^{MP} = \rho_{MP}u_{t-1}^{MP} + e_t^{MP}.
\]

\begin{resultbox}[title={Comparação: Regra CGG vs.\ Taylor (1993)}]
\begin{center}
\begin{tabular}{lll}
\toprule
 & Taylor (1993) & CGG (2000) \\
\midrule
Variáveis & Contemporâneas: $\pi_t$, $y_t$ & Expectacionais: $E_t[\pi_{t+1}]$, $E_t[y_{t+1}]$ \\
Estimável? & Sim (sem expectativas) & Exige GMM/expectativas \\
Período pré-Volcker & $\hat{\phi}_\pi \approx 0{,}83 < 1$ & Indeterminação (G.\ Inflação) \\
Período pós-Volcker & $\hat{\phi}_\pi \approx 2{,}15 > 1$ & Determinação (G.\ Moderação) \\
\bottomrule
\end{tabular}
\end{center}
\end{resultbox}

\subsection*{VI.2 Princípio de Taylor e Condição de Blanchard-Kahn}

\textbf{Sistema sem choques:}
\[
\begin{pmatrix}E_t[y_{t+1}]\\E_t[\pi_{t+1}]\end{pmatrix} = M\begin{pmatrix}y_t\\\pi_t\end{pmatrix}.
\]
Como $y_t$ e $\pi_t$ são variáveis de salto (forward-looking), Blanchard-Kahn exige \textbf{ambos os autovalores fora do círculo unitário}.

\begin{resultbox}[title={Princípio de Taylor e Determinação}]
\[
\phi_\pi > 1 \quad\text{(Princípio de Taylor)}
\]
Condição precisa: $\kappa(\phi_\pi-1) + (1-\beta)\phi_y > 0$.\\
Para $\beta\approx 1$: reduz-se a $\phi_\pi > 1$.
\end{resultbox}

\setcounter{section}{0}

% ==========================================================================
\section{Rigidez Nominal e Efeitos Reais da Moeda (Lista II Q1)}
% ==========================================================================

\subsection*{Enunciado}
Com base em modelos surgidos a partir da década de 70, a ocorrência de alguma forma de rigidez nominal na economia é essencial para: a) Efeitos nominais da política monetária; b) Efeitos reais de longo prazo de choques nos dispêndios; c) Coordenação entre políticas monetária e fiscal; d) Otimização dos agentes; \textbf{e) Efeitos reais de curto prazo dos choques monetários.}

\begin{provabox}[title={Resposta: alternativa (e)}]
A rigidez nominal --- preços ou salários que não ajustam instantaneamente --- é o mecanismo que permite que choques monetários tenham \textbf{efeitos reais no curto prazo} (impacto sobre produto e emprego). Sem rigidez nominal, qualquer choque monetário seria absorvido imediatamente por ajuste de preços, sem afetar variáveis reais. As opções (a)--(d) descrevem fenômenos que não dependem necessariamente de rigidez nominal.
\end{provabox}

% ==========================================================================
\section{Otimização Estática NK: $C^*$ e $L^*$ (Lista II Q2)}
% ==========================================================================

\subsection*{Enunciado}
Seja $U = C^{1-\theta} - \tfrac{1}{\gamma}L^\gamma$, $\gamma>1$, $\theta>0$, com restrição $C = WL$ ($W$ salário real). Calcule os níveis ótimos $C^*$ e $L^*$.

\subsection*{Resolução}

\textbf{Passo 1.} Lagrangiano ($\lambda$ é o multiplicador):
\[
\mathcal{L} = C^{1-\theta} - \frac{1}{\gamma}L^\gamma + \lambda(WL - C). \tag{1}
\]

\textbf{Passo 2.} CPO para $C$ ($\partial\mathcal{L}/\partial C = 0$):
\[
(1-\theta)C^{-\theta} - \lambda = 0 \;\Rightarrow\; \lambda = (1-\theta)C^{-\theta}. \tag{2}
\]

\textbf{Passo 3.} CPO para $L$ ($\partial\mathcal{L}/\partial L = 0$):
\[
-L^{\gamma-1} + \lambda W = 0 \;\Rightarrow\; \lambda = \frac{L^{\gamma-1}}{W}. \tag{3}
\]

\textbf{Passo 4.} Igualando (2) e (3):
\[
(1-\theta)C^{-\theta} = \frac{L^{\gamma-1}}{W}. \tag{4}
\]

\begin{resultbox}[title={Ótimos $C^*$ e $L^*$ (Q2)}]
\[
C^* = \left[\frac{W(1-\theta)}{L^{\gamma-1}}\right]^{1/\theta}, \qquad L^* = \bigl[W(1-\theta)\,C^{-\theta}\bigr]^{1/(\gamma-1)}.
\]
A condição (4) é análoga à condição $W/P = MRS$ entre consumo e trabalho: salário real = desutilidade marginal do trabalho / utilidade marginal do consumo.
\end{resultbox}

% ==========================================================================
\section{Concorrência Monopolística e Produto de Equilíbrio (Lista II Q3)}
% ==========================================================================

\subsection*{Enunciado}
Em concorrência monopolística o produto de equilíbrio é $Y = \left[\tfrac{\eta-1}{\eta}\right]^{1/(\gamma-1)}$, sendo $\eta>1$ a elasticidade da demanda e $\gamma>1$ a desutilidade do trabalho. Interprete os efeitos de $\eta$ e $\gamma$ sobre $Y$.

\subsection*{Derivação e Interpretação}

Da FOC da família ($MU_C = 1$, produção $Y=L$) e da condição de markup da firma ($P = \eta/(\eta-1)\cdot W$):
\[
\frac{W}{P} = \frac{\eta-1}{\eta} = L^{\gamma-1} = Y^{\gamma-1}
\;\Rightarrow\; Y = \left[\frac{\eta-1}{\eta}\right]^{1/(\gamma-1)}.
\]

\begin{resultbox}[title={Efeitos de $\eta$ e $\gamma$ sobre $Y$ (Q3)}]
\textbf{$\eta\uparrow$:} $(\eta-1)/\eta\uparrow$ $\Rightarrow$ menor markup (menos poder de mercado) $\Rightarrow$ salário real maior $\Rightarrow$ \textbf{$Y\uparrow$}.\\[3pt]
\textbf{$\gamma\uparrow$:} expoente $1/(\gamma-1)\downarrow$ (elasticidade da oferta de trabalho cai) $\Rightarrow$ trabalhadores respondem menos a incentivos salariais $\Rightarrow$ \textbf{$Y\downarrow$}.
\end{resultbox}

% ==========================================================================
\section{Nível de Preços de Equilíbrio (Lista II Q4)}
% ==========================================================================

\subsection*{Enunciado}
Com demanda agregada $Y = M/P$ e $Y$ do exercício anterior, deduza $P^*$. O que ocorre com $P$ quando $\gamma\uparrow$?

\subsection*{Resolução}

\[
Y = \frac{M}{P} \;\Longrightarrow\; P = \frac{M}{Y} = \frac{M}{\left[\frac{\eta-1}{\eta}\right]^{1/(\gamma-1)}}.
\]

\begin{resultbox}[title={Nível de Preços (Q4)}]
\[
\boxed{P^* = \frac{M}{\left[\frac{\eta-1}{\eta}\right]^{1/(\gamma-1)}}.}
\]
\textbf{$\gamma\uparrow$:} $Y\downarrow$ (Q3) $\Rightarrow$ dado $M$, excesso de demanda nominal $\Rightarrow$ \textbf{$P\uparrow$}. Maior desutilidade do trabalho contrai a oferta, reduz $Y$ e eleva preços.
\end{resultbox}

% ==========================================================================
\section{Modelo de Informação Incompleta de Lucas (Lista II Q5)}
% ==========================================================================

\subsection*{Enunciado}
No modelo de Lucas, a firma $i$ decide quanto produzir com base em $y_i = [1/(\gamma-1)]\,E[r_i|p_i]$, onde $r_i = p_i - p$ é o preço relativo estimado após a decisão. Explique flutuações de produto no curto prazo e como a rigidez de preços emerge.

\subsection*{Resolução}

A firma não conhece \emph{a priori} o nível geral de preços $p$ no momento da decisão de produção. Ela \textbf{extrai sinal} a partir de seus próprios preços passados e correntes $p_i$ para estimar $r_i = p_i - p$.

\textbf{Mecanismo de curto prazo:}
\begin{enumerate}[leftmargin=1.5em, itemsep=2pt]
  \item Choque positivo de demanda $\Rightarrow$ $p_i\uparrow$.
  \item Firma interpreta parte do aumento de $p_i$ como $r_i\uparrow$ (aumento de preço relativo).
  \item Responde elevando produção: $y_i = [1/(\gamma-1)]\,E[r_i|p_i]\uparrow$.
\end{enumerate}

\textbf{Rigidez de preços (intuitiva):} como a firma atribui parte do choque nominal a preço relativo, ela não eleva $p_i$ na mesma proporção do choque $\Rightarrow$ ajuste incompleto de preços $\Rightarrow$ rigidez informacional.

\begin{resultbox}[title={Lucas: Extração de Sinal (Q5)}]
\textbf{Curto prazo:} informação incompleta $\Rightarrow$ choque nominal confundido com real $\Rightarrow$ $y\uparrow$.\\
\textbf{Longo prazo:} $E[p] \to p$ (aprendizado) $\Rightarrow$ sem efeito real persistente.\\
\textbf{Elasticidade de produção:} $1/(\gamma-1)$. Quanto maior $\gamma$, menor a resposta produtiva a estímulos de preço relativo.
\end{resultbox}

% ==========================================================================
\section{Curva de Lucas e Neutralidade de Longo Prazo (Lista II Q6)}
% ==========================================================================

\subsection*{Enunciado}
Seja a Curva de Lucas $y = b(p - E[p])$, $b>0$. Interprete o produto no curto prazo após um choque monetário positivo e explique por que o produto converge ao nível normal no longo prazo.

\subsection*{Resolução}

\textbf{Curto prazo:} choque $m\uparrow$ (além de $E[m]$) $\Rightarrow$ $p > E[p]$ $\Rightarrow$ $y = b(p - E[p]) > 0$. Firmas confundem maior demanda nominal com maior preço relativo (Q5) e elevam produção.

\textbf{Longo prazo --- Neutralidade:} com expectativas racionais, os agentes aprendem a nova regra monetária ($E[m]\nearrow$, $E[p]\nearrow$):
\[
p = m = E[p] = E[m] \;\Rightarrow\; y = b(p - E[p]) = 0.
\]

\begin{resultbox}[title={Curva de Lucas: Curto e Longo Prazo (Q6)}]
\textbf{Curto prazo:} $m > E[m]$ $\Rightarrow$ $p > E[p]$ $\Rightarrow$ \textbf{$y > 0$} (efeito real via surpresa monetária).\\
\textbf{Longo prazo:} $\Delta E[m] > 0$ e $\Delta E[p] > 0$ $\Rightarrow$ $p - E[p] \to 0$ $\Rightarrow$ \textbf{$y \to 0$} (neutralidade).\\
Uma vez assimilada a nova regra de PM, ela deixa de ter efeitos reais.
\end{resultbox}

% ==========================================================================
\section{Crítica de Lucas e Mudança de Regra Monetária (Lista II Q7)}
% ==========================================================================

\subsection*{Enunciado}
Um BC usa séries temporais históricas para prever efeito expansionista de uma mudança de regra monetária. Interprete a eficácia dessa política segundo a Crítica de Lucas.

\subsection*{Resolução}

\begin{resultbox}[title={Crítica de Lucas (Q7)}]
As relações macroeconômicas estimadas com dados históricos são \textbf{condicionadas ao estado de expectativas vigente naquele período}. Uma mudança estrutural de regra monetária \textbf{altera as expectativas} dos agentes, modificando as próprias relações macroeconômicas.\\[3pt]
\textbf{Consequência:} previsões baseadas em séries passadas perdem validade após a mudança de regra.\\[3pt]
\textbf{Condição de validade das previsões:} a mudança de PM não pode ser percebida como estrutural (permanente) pelos agentes.
\end{resultbox}

\begin{provabox}[title={Implicação Metodológica da Crítica de Lucas (Q7)}]
A Crítica motivou a micro-fundamentação dos modelos DSGE: parâmetros estruturais (preferências, tecnologia) são invariantes a mudanças de política; relações reduzidas estimadas não o são. O modelo NK canônico CGG (Q11) é uma resposta direta a essa crítica.
\end{provabox}

% ==========================================================================
\section{Papel de $\phi$ nos Modelos de Fischer e Taylor (Lista II Q8)}
% ==========================================================================

\subsection*{Enunciado}
Como o grau de rigidez de preços (a partir de $\phi$) afeta os efeitos reais da política monetária em: a) Modelo de Fischer; b) Modelo de Taylor?

\subsection*{Resolução}

\textbf{(Ver derivações completas em Parte V, Seções V.2--V.3.)}

\textbf{a) Modelo de Fischer:}
\[
p_t = E_{t-2}m_t + \frac{\phi}{1+\phi}\left[E_{t-1}m_t - E_{t-2}m_t\right], \qquad
y_t = \frac{1}{1+\phi}\left[E_{t-1}m_t - E_{t-2}m_t\right] + \left[m_t - E_{t-1}m_t\right].
\]
$\phi$ é o \textbf{inverso do grau de rigidez real}: $\phi\uparrow$ $\Rightarrow$ revisões de expectativas têm maior efeito sobre \textbf{preços} e \textbf{menor} sobre o \textbf{produto}.

\textbf{b) Modelo de Taylor:}
\[
y_t = \lambda y_{t-1} + \frac{1+\lambda}{2}u_t, \qquad \lambda = \frac{1-\phi}{1+\phi}.
\]
Relação inversa entre $\lambda$ e $\phi$: $\phi\uparrow$ $\Rightarrow$ $\lambda\downarrow$ $\Rightarrow$ menos persistência e menor efeito de choques de demanda $u_t$.

\begin{resultbox}[title={Fischer vs.\ Taylor: Papel de $\phi$ (Q8)}]
\begin{center}
\begin{tabular}{lll}
\toprule
Modelo & $\phi\uparrow$ (menor rigidez real) & Resultado \\
\midrule
Fischer & Revisões $[E_{t-1}m_t - E_{t-2}m_t]\uparrow$ no preço & Menor efeito no produto \\
Taylor & $\lambda = (1-\phi)/(1+\phi)\downarrow$ & Menos persistência de $y_t$ \\
\bottomrule
\end{tabular}
\end{center}
\textbf{Caso especial $\phi = 1$:} no Taylor, $\lambda=0$ $\Rightarrow$ moeda neutra; no Fischer, inovação ainda tem efeito real pleno.
\end{resultbox}

% ==========================================================================
\section{Calvo: Papel de $\theta$ e $\alpha$ na CPNK (Lista II Q9)}
% ==========================================================================

\subsection*{Enunciado}
Na CPNK $\pi_t = \kappa y_t + \beta E_t\pi_{t+1}$ com $\kappa = \frac{\alpha[1-(1-\alpha)\beta]\phi}{1-\alpha}$ e $\phi = \theta+\gamma-1$: a) papel de $\theta$; b) papel de $\alpha$.

\subsection*{Resolução}

\textbf{(Ver derivação completa em Parte V, Seção V.4.)}

\begin{resultbox}[title={Papel de $\theta$ e $\alpha$ na CPNK (Q9)}]
\textbf{a) Aversão ao risco $\theta$:} $\phi = \theta+\gamma-1$. $\theta\uparrow$ $\Rightarrow$ $\phi\uparrow$ (menor rigidez real) $\Rightarrow$ $\kappa\uparrow$. Um aumento de $\theta$ \textbf{aumenta a sensibilidade da inflação ao hiato do produto}.\\[4pt]
\textbf{b) Fração de ajuste $\alpha$:} quanto maior $\alpha$, mais firmas ajustam preços em cada período $\Rightarrow$ $\kappa\uparrow$. Intuitivamente: mais flexibilidade de preços $\Rightarrow$ maior transmissão do produto para a inflação.
\end{resultbox}

% ==========================================================================
\section{Calvo vs.\ Christiano-Eichenbaum-Evans (2005) (Lista II Q10)}
% ==========================================================================

\subsection*{Enunciado}
Explique a diferença entre Calvo (1983) e Christiano et al.\ (2005) e como o segundo permite desinflação com custo positivo.

\subsection*{Resolução}

\textbf{(Ver derivação completa em Parte V, Seção V.5.)}

\textbf{Calvo puro:} $\pi_t = \kappa y_t + \beta E_t\pi_{t+1}$. Desinflação crível ($E_t\pi_{t+1}\downarrow$) implica $y\geq 0$ --- contradiz a evidência (Ball 1994b: produto abaixo do normal em 27/28 episódios de desinflação).

\textbf{CEE (2005):} contratos com \textbf{indexação} repassam automaticamente a inflação passada $\Rightarrow$ CPNK híbrida:
\[
\pi_t = \gamma\pi_{t-1} + (1-\gamma)E_t\pi_{t+1} + \chi y_t, \quad 0\leq\gamma\leq 1.
\]

\begin{resultbox}[title={Custo de Desinflação no CEE (Q10)}]
Condição: $\gamma > 1/2$ (componente backward-looking domina).\\
Com $\pi_{t-1} > \pi_t > E_t\pi_{t+1}$ (trajetória declinante), satisfazer a CPNK exige $y_t < 0$ (\textbf{recessão}).\\[3pt]
\textbf{Exemplo numérico:} $\gamma=0{,}8$, $\chi=0{,}9$, $\pi_{t-1}=6\%$, $E_t\pi_{t+1}=1{,}5\%$, $\pi_t=3{,}5\%$:\\
$y_t = [3{,}5 - 0{,}8\times 6 - 0{,}2\times 1{,}5]/0{,}9 = -1{,}6/0{,}9 \approx -1{,}7\%$.
\end{resultbox}

% ==========================================================================
\section{Modelo NK Canônico CGG: $\theta$ e Choques de Oferta (Lista II Q11)}
% ==========================================================================

\subsection*{Enunciado}
Sob o modelo NK canônico (IS + CPNK + Regra CGG): i) como $\theta$ impacta as relações estruturais? ii) como choques negativos de oferta influenciam a condução da PM?

\subsection*{Resolução}

\textbf{(Ver modelo completo em Parte VI, Seção VI.1.)}

\[
y_t = E_t[y_{t+1}] - \frac{1}{\theta}r_t + u_t^{IS}, \quad
\pi_t = \beta E_t[\pi_{t+1}] + \kappa y_t + u_t^\pi, \quad
r_t = \phi_\pi E_t[\pi_{t+1}] + \phi_y E_t[y_{t+1}] + u_t^{MP}.
\]

\begin{resultbox}[title={i) Papel de $\theta$ no Modelo CGG (Q11)}]
\textbf{Curva IS:} $\theta\uparrow$ $\Rightarrow$ $1/\theta\downarrow$ $\Rightarrow$ menor eficácia da PM sobre o hiato ($y$ menos sensível a $r$). Para neutralizar um choque de demanda, o BC precisa elevar $r$ mais.\\[3pt]
\textbf{CPNK:} $\theta\uparrow$ $\Rightarrow$ $\phi = \theta+\gamma-1\uparrow$ $\Rightarrow$ $\kappa\uparrow$ $\Rightarrow$ maior impacto inflacionário de choques $\Rightarrow$ PM deve reagir mais agressivamente.\\[3pt]
\textbf{Regra CGG:} quanto maiores $\theta$ e $\kappa$ (e menor $\beta$), maiores os coeficientes ótimos $\phi_\pi$ e $\phi_y$.
\end{resultbox}

\begin{resultbox}[title={ii) Choques Negativos de Oferta (Q11)}]
Choque $u_t^\pi > 0$ na CPNK: \textbf{estagflação} ($\pi\uparrow$ e $y\downarrow$ simultaneamente).\\
\textbf{Dilema do BC:} estabilizar $\pi$ exige elevar $r$ (aprofunda $y\downarrow$); estabilizar $y$ exige acomodar $\pi$.\\
\textbf{Persistência alta} ($\rho_\pi$ elevado): o BC pode preferir desinflação custosa a arriscar desancoragem das expectativas e sobreposição de novos choques.\\
\textbf{Persistência baixa}: acomodar transitoriamente (bandas de tolerância) pode ser ótimo --- choques tendem a se reverter.
\end{resultbox}

\begin{provabox}[title={Rigidez Nominal e Efeitos Reais da Meda --- Para a Prova}]
A progressão lógica dos modelos NK é: custos de menu/informação incompleta (Lucas) $\to$ contratos escalonados (Fischer/Taylor) $\to$ probabilidade de Calvo $\to$ indexação (CEE) $\to$ modelo canônico (CGG). Cada passo adiciona um mecanismo de transmissão de choques monetários e de inércia inflacionária.
\end{provabox}


% ==========================================================================
\appendix
% ==========================================================================

\newpage
\section*{\large\color{cyanrbc} Apêndice A --- Questões-tipo (Fischer, Taylor, Calvo, CEE e CGG)}
\addcontentsline{toc}{section}{Apêndice A --- Questões-tipo com novos modelos}

\subsection*{Q-tipo 1 --- Fischer vs.\ Taylor: preços predeterminados vs.\ fixos}

\begin{enumerate}[label=(\alph*)]
  \item No Fischer, os preços são completamente rígidos (não podem mudar intra-contrato).
  \item No Taylor, os preços são apenas predeterminados: podem ser ajustados se houver nova informação.
  \item No Fischer, o escalonamento permite que inovações monetárias em $t-1$ afetem $y_t$; no Taylor, isso também ocorre via persistência $\lambda$.
  \item Ambos os modelos implicam neutralidade da moeda para qualquer $\phi > 0$.
\end{enumerate}

\begin{resultbox}[title={Gabarito: (c)}]
\textbf{Correta: (c).} Fischer tem preços predeterminados (mas a coorte que fixa em $t-1$ pode usar informação de $t-1$). Taylor tem preços fixos por dois períodos. Em ambos, o escalonamento gera efeitos reais de surpresas monetárias. (a) e (b) trocam as definições. (d) FALSO: ambos implicam não-neutralidade da moeda.
\end{resultbox}

\subsection*{Q-tipo 2 --- Rigidez real no modelo de Taylor: efeito de $\phi\downarrow$}

Se $\phi$ decresce de 1 para 0,5 no modelo de Taylor:

\begin{enumerate}[label=(\alph*)]
  \item $\lambda$ decresce: menos persistência do produto.
  \item $\lambda$ aumenta: maior persistência do produto.
  \item O produto torna-se completamente inelástico a choques monetários.
  \item Apenas a inflação é afetada; o produto permanece inalterado.
\end{enumerate}

\begin{resultbox}[title={Gabarito: (b)}]
\textbf{Correta: (b).} $\lambda = (1-\phi)/(1+\phi)$. Para $\phi=1$: $\lambda=0$. Para $\phi=0{,}5$: $\lambda = 0{,}5/1{,}5 \approx 0{,}33 > 0$. Maior rigidez real ($\phi\downarrow$) $\Rightarrow$ $\lambda\uparrow$ $\Rightarrow$ mais persistência.
\end{resultbox}

\subsection*{Q-tipo 3 --- CPNK Calvo: notação do professor}

Na versão do professor (slides), a inclinação da CPNK é $\kappa = \alpha[1-(1-\alpha)\beta]\phi/(1-\alpha)$. Um aumento em $\phi$ (mantendo $\alpha$ e $\beta$ fixos):

\begin{enumerate}[label=(\alph*)]
  \item Reduz $\kappa$: menor rigidez real $\Rightarrow$ CPNK mais plana.
  \item Aumenta $\kappa$: menor rigidez real $\Rightarrow$ CPNK mais íngreme.
  \item Não afeta $\kappa$: $\phi$ não aparece na fórmula.
  \item Reduz $\beta$: o desconto intertemporal cai.
\end{enumerate}

\begin{resultbox}[title={Gabarito: (b)}]
\textbf{Correta: (b).} $\kappa = \alpha[1-(1-\alpha)\beta]\phi/(1-\alpha)$. $\phi$ entra linearmente no numerador: $\phi\uparrow$ $\Rightarrow$ $\kappa\uparrow$. Economicamente: menor rigidez real (salário real mais sensível ao produto) $\Rightarrow$ pass-through maior de $y$ para $\pi$ $\Rightarrow$ CPNK mais íngreme.
\end{resultbox}

\subsection*{Q-tipo 4 --- CEE vs.\ Calvo puro: custo de sacrifício}

Para que a CPNK híbrida de CEE ($\pi_t = \gamma\pi_{t-1} + (1-\gamma)E_t\pi_{t+1} + \chi y_t$) implique custo positivo de desinflação:

\begin{enumerate}[label=(\alph*)]
  \item $\gamma < 1/2$: componente forward-looking domina $\Rightarrow$ custo positivo.
  \item $\gamma > 1/2$: componente backward-looking domina $\Rightarrow$ custo positivo.
  \item $\gamma = 1/2$: custo máximo de desinflação.
  \item $\chi > 0$ é condição suficiente para custo positivo, independente de $\gamma$.
\end{enumerate}

\begin{resultbox}[title={Gabarito: (b)}]
\textbf{Correta: (b).} Para desinflação com custo: quando $E_t\pi_{t+1} < \pi_t$ (tendência declinante), precisamos $y_t < 0$ para satisfazer a CPNK híbrida. Isso requer que $\gamma\pi_{t-1}$ pese mais que $(1-\gamma)E_t\pi_{t+1}$ no lado direito, i.e., $\gamma > 1-\gamma$ $\Leftrightarrow$ $\gamma > 1/2$.
\end{resultbox}

\subsection*{Q-tipo 5 --- Regra CGG vs.\ Taylor simples}

A diferença entre a regra CGG $r_t = \phi_\pi E_t[\pi_{t+1}] + \phi_y E_t[y_{t+1}]$ e Taylor simples $r_t = \phi_\pi\pi_t + \phi_y y_t$ é:

\begin{enumerate}[label=(\alph*)]
  \item CGG é mais simples pois não exige previsões.
  \item CGG reage a expectativas futuras; Taylor reage a valores contemporâneos.
  \item As duas regras são equivalentes se os agentes têm expectativas racionais.
  \item CGG implica sempre $\phi_\pi < 1$ (política passiva obrigatória).
\end{enumerate}

\begin{resultbox}[title={Gabarito: (b)}]
\textbf{Correta: (b).} CGG usa $E_t[\pi_{t+1}]$ e $E_t[y_{t+1}]$ (forward-looking); Taylor (1993) usa $\pi_t$ e $y_t$ (contemporânea). A distinção é empiricamente importante: CGG requer estimação por GMM. (c) FALSO: mesmo com expectativas racionais, $E_t[\pi_{t+1}] \neq \pi_t$ em geral.
\end{resultbox}

\subsection*{Q-tipo 6 --- Ball (1994b): evidência e inércia inflacionária}

Ball (1994b) analisou 28 episódios de desinflação em 9 países e encontrou que:

\begin{enumerate}[label=(\alph*)]
  \item Em 28 de 28 casos, o produto ficou acima do normal durante a desinflação (boom desinflacionário).
  \item Em 27 de 28 casos, o produto ficou abaixo do normal (recessão).
  \item Em todos os casos, a desinflação foi crível e sem custo real.
  \item Os episódios mostraram que a CPNK pura é consistente com os dados.
\end{enumerate}

\begin{resultbox}[title={Gabarito: (b)}]
\textbf{Correta: (b).} Ball (1994b): 27/28 episódios com produto abaixo do normal. Isso refuta a CPNK pura de Calvo (que implica desinflação crível sem custo) e motiva modelos com inércia real como o CEE. (a), (c) e (d) são todos incorretos.
\end{resultbox}

\subsection*{Q-tipo 7 --- RCK: Questões adicionais da Lista I}

Sobre a regra de Keynes-Ramsey $\dot{c}/c = [f'(k)-\delta-\rho-\theta g]/\theta$, qual afirmação é correta?

\begin{enumerate}[label=(\alph*)]
  \item $\rho\uparrow$ eleva o crescimento do consumo.
  \item $\theta\uparrow$ (maior aversão ao risco) eleva $\dot{c}/c$.
  \item $n\uparrow$ aparece em $\dot{c}/c$ com sinal positivo.
  \item $\rho\uparrow$ afeta \emph{negativamente} $\dot{c}/c$.
\end{enumerate}

\begin{resultbox}[title={Gabarito: (d)}]
\textbf{Correta: (d).} $\rho$ entra com sinal $-$ no numerador: maior impaciência $\Rightarrow$ menor crescimento do consumo. (a) inverte o sinal. (b) FALSO: $\theta\uparrow$ reduz $1/\theta$ e pode elevar $r^*$, mas o efeito líquido sobre $\dot{c}/c$ é redução (via EE). (c) FALSO: $n$ não aparece em $\dot{c}/c$.
\end{resultbox}

% ==========================================================================
\section*{\large\color{cyanrbc} Apêndice B --- Tabelas de Calibração Unificadas}
\addcontentsline{toc}{section}{Apêndice B --- Tabelas de Calibração Unificadas}
% ==========================================================================

\subsection*{B.1 Parâmetros dos Modelos de Crescimento}

\begin{center}
\begin{tabular}{llccc}
\toprule
Parâmetro & Descrição & Solow & RCK & RBC \\
\midrule
$\alpha$ & Participação do capital & 0,33 & 0,33 & 0,33 \\
$\delta$ & Taxa de depreciação & 0,025 & 0,025 & 0,025 \\
$n$ & Crescimento pop.\ & 0,01 & 0,01 & 0 \\
$g$ & Prog.\ tecnológico & 0,02 & 0,02 & 0 \\
$s$ & Taxa de poupança & 0,20 & --- (end.) & --- \\
$\rho$ & Taxa de desconto (cont.) & --- & 0,04 & --- \\
$\theta$ & Aversão ao risco (CRRA) & --- & 2 & 1 (log) \\
$\beta$ & Fator de desconto (disc.) & --- & --- & 0,99 \\
$b$ & Peso do lazer & --- & --- & $\approx 1{,}55$ \\
$\rho_z$ & Persistência choque PTF & --- & --- & 0,95 \\
\bottomrule
\end{tabular}
\end{center}

\subsection*{B.2 Parâmetros do Modelo NK}

\begin{center}
\begin{tabular}{llcc}
\toprule
Parâmetro & Descrição & Valor base & Referência \\
\midrule
$\beta$ & Fator de desconto (trimestral) & $0{,}99$ & Padrão \\
$\theta$ & Inv.\ EIS (= $\sigma$ em Galí) & $1{,}0$ & Log-utilidade \\
$\gamma$ & Curvatura desut.\ trabalho & $2{,}0$ & $\Rightarrow \phi = \theta+\gamma-1 = 2$ \\
$\phi$ & Rigidez real ($=\theta+\gamma-1$) & $2{,}0$ & Slides prof. \\
$\alpha_{Calvo}$ & Prob.\ não-ajuste & $0{,}75$ & Galí (2008) \\
$\kappa$ & Slope CPNK & $\approx 0{,}1$ & Calibrado \\
$\phi_\pi$ & Peso inflação (CGG) & $1{,}5$ & Taylor (1993) \\
$\phi_y$ & Peso produto (CGG) & $0{,}5$ & Taylor (1993) \\
$\rho_{IS}$ & Persist.\ choque IS & $0{,}5$ & Estimado \\
$\rho_\pi$ & Persist.\ choque custo-push & $0{,}5$ & Estimado \\
$\gamma_{CEE}$ & Peso backward (híbrida) & $0{,}5$ & CEE (2005) \\
\bottomrule
\end{tabular}
\end{center}

\subsection*{B.3 Slope da CPNK ($\kappa$) como Função de $\alpha_{Calvo}$ e $\phi$}

\begin{center}
\begin{tabular}{ccccc}
\toprule
$\alpha_{Calvo}$ & Duração esperada & $1-(1-\alpha_{Calvo})\beta$ & $\kappa$ (com $\phi=2$) & $\kappa$ (com $\phi=1$) \\
\midrule
0,50 & 2 tri. & $0{,}505$ & $\approx 1{,}010$ & $\approx 0{,}505$ \\
0,60 & 2,5 tri. & $0{,}406$ & $\approx 0{,}541$ & $\approx 0{,}271$ \\
0,75 & 4 tri. & $0{,}257$ & $\approx 0{,}171$ & $\approx 0{,}086$ \\
0,90 & 10 tri. & $0{,}109$ & $\approx 0{,}024$ & $\approx 0{,}012$ \\
\bottomrule
\end{tabular}
\end{center}
Usando $\kappa = \alpha_{Calvo}[1-(1-\alpha_{Calvo})\beta]\phi/(1-\alpha_{Calvo})$.

\subsection*{B.4 Momentos do Modelo RBC vs.\ Dados (EUA Trimestrais)}

\begin{center}
\begin{tabular}{lccc}
\toprule
Variável & Dados EUA & RBC & Interpretação \\
\midrule
$\sigma(Y)$ & 1,72\% & 1,72\% & Calibrado pelo prof. \\
$\sigma(C)/\sigma(Y)$ & 0,74 & 0,48 & RBC subestima (canal 3) \\
$\sigma(I)/\sigma(Y)$ & 2,98 & 3,15 & Consistente (canal 2) \\
$\sigma(N)/\sigma(Y)$ & 0,95 & 0,77 & Subestima (trabalho intensivo) \\
$\text{corr}(C,Y)$ & 0,88 & 0,94 & Supercorrelação do consumo \\
$\text{corr}(I,Y)$ & 0,80 & 0,99 & Investimento muito correlacionado \\
\bottomrule
\end{tabular}
\end{center}

\subsection*{B.5 Conexão entre Modelos}

\begin{center}
\begin{formulabox}[title={Lógica da Sequência de Modelos}]
\begin{enumerate}[leftmargin=1.5em, itemsep=3pt]
  \item \textbf{Solow}: crescimento exógeno ($s$ exógena). Ponto de partida.
  \item \textbf{RCK}: endogeneíza $s$ via otimização intertemporal. Adiciona Keynes-Ramsey.
  \item \textbf{Diamond}: gerações sobrepostas. Possível ineficiência dinâmica.
  \item \textbf{RBC}: adiciona estocasticidade (PTF). Mercados competitivos + preços flexíveis.
  \item \textbf{Fischer/Taylor}: rigidez nominal via contratos escalonados. Moeda não-neutra.
  \item \textbf{Calvo}: rigidez estocástica (probabilidade $\alpha$ de não ajustar). Gera CPNK.
  \item \textbf{CEE}: adiciona indexação à Calvo. Gera CPNK híbrida com inércia real.
  \item \textbf{CGG}: modelo de três equações completo (IS + CPNK + Regra forward-looking).
\end{enumerate}
\end{formulabox}
\end{center}

\begin{center}
\rule{0.5\linewidth}{0.4pt}\\[0.5em]
{\small Guia unificado para a prova única de Macroeconomia I --- PPGEco-UFES 2026.\\
Código Python, soluções numéricas e material complementar: \texttt{github.com/fcarva/macro}}
\end{center}

\end{document}
